\documentclass[aspectratio=169,11pt]{beamer}

\usetheme{Madrid}
\usecolortheme{default}
\usefonttheme{professionalfonts}
\setbeamertemplate{navigation symbols}{}
\setbeamertemplate{footline}[frame number]

\usepackage{amsmath, amssymb, booktabs, array}

\title{Experiments, Measurement, and Behavioral Identification in Labor}
\subtitle{Behavioral Labor -- Week 7}
\author{Teaching deck draft}
\date{}

\begin{document}

\begin{frame}
  \titlepage
\end{frame}

\begin{frame}{Central question}
\begin{itemize}
    \item How do behavioral labor economists move from a behavioral object to an actual empirical method?
    \item Week 7 is the toolkit lecture: measurement, experiments, hazards, bunching, quasi-experiments, and structural estimation.
    \item The practical arc: object, labor margin, variation, estimand, method, interpretation.
    \item The recurring warning: a setting is not an estimand.
\end{itemize}
\end{frame}

\begin{frame}{Bridge from Weeks 1--6}
\small
\begin{tabular}{p{0.25\linewidth} p{0.60\linewidth}}
\toprule
Earlier object & Week 7 empirical task \\
\midrule
Preferences & recover dynamic, reference-dependent, or social preference parameters \\
Beliefs & elicit, validate, and link expectations to search and choices \\
Attention & separate salience, knowledge, and complexity from true incentives \\
Workplace behavior & estimate treatment effects and recover mechanisms \\
Identity and norms & distinguish measurement, sorting, treatment, and transmission \\
\bottomrule
\end{tabular}

\vspace{0.7em}
Week 8 will apply this toolkit to firm, market, and equilibrium response.
\end{frame}

\begin{frame}{What counts as behavioral identification?}
\begin{itemize}
    \item Name the behavioral object: belief, preference, attention cost, reciprocity, norm sensitivity, or learning rule.
    \item Name the labor margin: search, exit, effort, take-up, applications, training, retention, or earnings.
    \item Name the variation: treatment, timing, threshold, schedule, exposure, survey panel, or model moments.
    \item Name the estimand: ITT, ATE, hazard effect, local elasticity, belief updating, or structural parameter.
    \item Name the strongest standard alternative.
\end{itemize}
\end{frame}

\begin{frame}{Measurement: latent object and noisy proxy}
\[
z_{ij} = \theta_i + \eta_{ij}
\]

\small
\begin{tabular}{p{0.22\linewidth} p{0.64\linewidth}}
\toprule
Object & Interpretation \\
\midrule
$\theta_i$ & latent behavioral object: belief, preference, attention, norm sensitivity \\
$z_{ij}$ & survey response, elicitation, task behavior, admin proxy, or model moment \\
$\eta_{ij}$ & measurement noise, demand effects, misunderstanding, timing mismatch \\
\bottomrule
\end{tabular}

\vspace{0.7em}
Behavioral labor often starts by asking whether the proxy really measures the object.
\end{frame}

\begin{frame}{Behavioral object to data design}
\scriptsize
\begin{tabular}{p{0.20\linewidth} p{0.22\linewidth} p{0.25\linewidth} p{0.20\linewidth}}
\toprule
Object & Labor setting & Data object & Main challenge \\
\midrule
Present bias & training, savings, effort & repeated choices, contract menus & preferences versus liquidity \\
Beliefs & job search & elicited expectations, linked outcomes & beliefs versus opportunities \\
Attention & schedules, take-up & admin records, notices, claims & salience versus constraints \\
Social preferences & workplace effort & real-effort output, gifts, wages & reciprocity versus selection \\
Norms & allocation, teams & coworkers, managers, surveys & treatment versus sorting \\
\bottomrule
\end{tabular}
\end{frame}

\begin{frame}{Field experiments in labor settings}
\[
Y_i = \alpha + \beta T_i + X_i'\gamma + \varepsilon_i
\]

\begin{itemize}
    \item Assignment identifies the treatment effect; OLS or ANCOVA is often the practical estimator.
    \item Baseline outcomes and strata improve precision.
    \item Randomization inference helps with small, clustered, or design-sensitive experiments.
    \item Heterogeneity matters when treatment works through prior beliefs, knowledge, constraints, or worker type.
    \item Mechanism interpretation still requires measurement.
\end{itemize}
\end{frame}

\begin{frame}{Information interventions}
\small
\begin{tabular}{p{0.28\linewidth} p{0.28\linewidth} p{0.28\linewidth}}
\toprule
Information object & Labor example & Interpretation issue \\
\midrule
Belief level & job-finding probabilities & did expectations update? \\
Salience & EITC or benefit reminders & knowledge or attention? \\
Procedural knowledge & how to claim or apply & information or hassle reduction? \\
Confidence & perceived eligibility & optimism or actual opportunity? \\
Learning speed & schedule understanding & temporary or persistent response? \\
\bottomrule
\end{tabular}

\vspace{0.7em}
The regression may be simple; the behavioral claim is not.
\end{frame}

\begin{frame}{Survey expectations and panels}
\begin{itemize}
    \item Repeated elicitation measures beliefs, bias, persistence, and updating.
    \item Panel fixed effects separate stable optimism or pessimism from within-person updating.
    \item Validation links beliefs to realized exits, wages, applications, or search effort.
    \item Forecast errors help distinguish beliefs from labor-market opportunities.
    \item Measurement error is central, not cosmetic.
\end{itemize}
\end{frame}

\begin{frame}{Job-search durations and hazards}
\[
h_{it} = \Pr(U_{i,t+1}=0 \mid U_{it}=1,\, X_{it}, B_{it})
\]

\begin{itemize}
    \item Search outcomes are often spells, exits, and censoring problems.
    \item Discrete-time logit or probit hazards fit period-by-period exit data.
    \item Cox-style hazards fit continuous duration settings.
    \item Competing risks matter when exits include jobs, training, inactivity, or benefit transitions.
    \item Behavioral question: beliefs, discouragement, learning, benefits, or opportunity?
\end{itemize}
\end{frame}

\begin{frame}{Bunching, salience, and nonlinear schedules}
\[
\hat e = \frac{\Delta b / b}{\Delta (1-\tau)/(1-\tau)}
\]

\begin{itemize}
    \item Nonlinear tax, transfer, earnings, and bonus schedules create local variation.
    \item Bunching, kink, and notch methods estimate local response around schedule features.
    \item Low bunching may mean low elasticities, adjustment costs, inattention, or poor knowledge.
    \item Knowledge and salience data discipline behavioral interpretation.
\end{itemize}
\end{frame}

\begin{frame}{Quasi-experiments in behavioral settings}
\small
\begin{tabular}{p{0.25\linewidth} p{0.28\linewidth} p{0.30\linewidth}}
\toprule
Variation & Method family & Behavioral use \\
\midrule
threshold or score & RD, kink, notch & eligibility, schedule knowledge, local response \\
staggered rollout & DiD, event study & information release, simplification, exposure \\
encouragement & IV & take-up when compliance is endogenous \\
policy timing & event study & learning, anticipation, dynamic adjustment \\
schedule reform & bunching, local elasticity & salience and incentive response \\
\bottomrule
\end{tabular}
\end{frame}

\begin{frame}{Structural estimation in behavioral labor}
\[
\hat \psi \in \arg\min_{\psi \in \Psi}
\left[m^{data}-m^{model}(\psi)\right]'W
\left[m^{data}-m^{model}(\psi)\right]
\]

\begin{itemize}
    \item $\psi$ can include present bias, reciprocity, search costs, learning rules, or attention costs.
    \item MLE, SMM, dynamic discrete choice, and calibrated moments are common tools.
    \item Structural work is useful for parameter recovery and counterfactuals.
    \item The price is model dependence: state which moments identify which parameters.
\end{itemize}
\end{frame}

\begin{frame}{Practical econometric methods cheat sheet}
\scriptsize
\begin{tabular}{p{0.30\linewidth} p{0.30\linewidth} p{0.25\linewidth}}
\toprule
Empirical setting & Common method & Add-ons \\
\midrule
Randomized labor intervention & OLS, ANCOVA & randomization inference, heterogeneity \\
Repeated beliefs panel & fixed effects, validation & forecast errors, measurement error \\
Job-search duration & logit/probit hazards, Cox & competing risks, dynamic controls \\
Nonlinear schedule & bunching, local elasticity & knowledge interactions \\
Policy or information shock & IV, RD, DiD, event study & spillovers, dynamics \\
Parameter recovery & MLE, SMM, dynamic model & fit, sensitivity \\
\bottomrule
\end{tabular}
\end{frame}

\begin{frame}{Heterogeneity is a method choice}
\begin{itemize}
    \item Behavioral effects often vary by baseline beliefs, knowledge, attention, liquidity, identity, or workplace context.
    \item Interactions are transparent when the moderator is pre-treatment.
    \item Random coefficients or hierarchical models can summarize richer heterogeneity.
    \item Machine-learning heterogeneity tools can be useful, but the behavioral object still needs interpretation.
    \item Never let heterogeneity replace a mechanism.
\end{itemize}
\end{frame}

\begin{frame}{Common failure modes}
\small
\begin{tabular}{p{0.30\linewidth} p{0.48\linewidth}}
\toprule
Failure & Diagnostic question \\
\midrule
Setting without estimand & what parameter is estimated? \\
Beliefs versus opportunities & are expectations validated against outcomes? \\
Treatment versus measurement & what information object actually moved? \\
Bunching overreach & preferences, knowledge, or adjustment costs? \\
Structural overreach & which moments identify which parameters? \\
Welfare over-claiming & behavior only, or welfare too? \\
\bottomrule
\end{tabular}
\end{frame}

\begin{frame}{Week 7 lab}
\begin{itemize}
    \item Reproduce: synthetic job-search information experiment anchored on Altmann--Falk--Jaeger--Zimmermann.
    \item Diagnose: behavioral object, margin, variation, method, assumptions, interpretation.
    \item Transfer: workplace gift exchange and social-preference recovery.
    \item Challenge: nonlinear schedule response or self-control at work.
\end{itemize}
\end{frame}

\begin{frame}{Bridge to Week 8}
\begin{itemize}
    \item Week 7 supplies the empirical toolkit.
    \item Week 8 asks what survives once worker-level behavioral wedges enter firms, markets, search, and sorting.
    \item The method discipline separates worker response, firm strategy, and market outcomes.
    \item Week 9 then turns that equilibrium view into behavioral public policy design.
\end{itemize}
\end{frame}

\begin{frame}{Takeaways}
\begin{itemize}
    \item Behavioral labor combines standard labor methods with behavioral measurement and modeling.
    \item Experiments estimate treatment effects; mechanisms require measurement.
    \item Beliefs panels call for validation and fixed effects.
    \item Search durations call for hazards; nonlinear schedules call for bunching and local elasticities.
    \item Structural tools are for parameter recovery and counterfactuals, not decorative language.
\end{itemize}
\end{frame}

\end{document}
